\documentclass[11pt]{article}

% NeurIPS/ICML style packages
\usepackage[utf8]{inputenc}
\usepackage[T1]{fontenc}
\usepackage{hyperref}
\usepackage{url}
\usepackage{booktabs}
\usepackage{amsfonts}
\usepackage{amsmath}
\usepackage{graphicx}
\usepackage{xcolor}
\usepackage{microtype}
\usepackage{natbib}
\usepackage{algorithm}
\usepackage{algorithmic}
\usepackage{listings}
\usepackage{subcaption}

% Page layout
\usepackage[margin=1in]{geometry}

% Code listing style
\lstset{
  basicstyle=\ttfamily\small,
  breaklines=true,
  frame=single,
  numbers=left,
  numberstyle=\tiny\color{gray},
  keywordstyle=\color{blue},
  commentstyle=\color{green!60!black},
  stringstyle=\color{red!80!black},
}

\title{Cloud-Gym: A Benchmark for LLM-Based Infrastructure-as-Code Repair}

\author{
  Author Name \\
  Affiliation \\
  \texttt{email@example.com}
}

\date{}

\begin{document}

\maketitle

\begin{abstract}
Infrastructure-as-Code (IaC) configurations are critical to modern cloud deployments, yet LLM-based repair of broken IaC remains underexplored due to the lack of standardized benchmarks. We present \textbf{Cloud-Gym}, a benchmark and training data generation framework for evaluating LLM repair of Terraform and CloudFormation configurations. Cloud-Gym features a 37-fault taxonomy spanning 8 categories (syntactic, reference, semantic, dependency, provider, security, cross-resource, and intrinsic function faults), a programmatic inversion engine that generates validated broken configurations from gold instances, and an evaluation harness using the pass@$k$ metric. We collect gold configurations from GitHub, the Terraform Registry, and AWS sample repositories, generate thousands of training pairs via both rule-based and LLM-assisted fault injection, and establish baseline repair rates using open-weight models. Our results show that while models handle syntactic faults reasonably well, semantic and cross-resource faults remain challenging, motivating further research in IaC-specific code repair.
\end{abstract}

\section{Introduction}
\label{sec:introduction}

Infrastructure-as-Code (IaC) has become the standard for managing cloud resources, with tools like Terraform and CloudFormation enabling declarative specification of complex infrastructure. As IaC configurations grow in complexity, misconfigurations become a significant source of deployment failures, security vulnerabilities, and operational incidents.

Recent advances in large language models (LLMs) have shown promise for automated code repair across general-purpose programming languages. However, IaC repair presents unique challenges:
\begin{itemize}
    \item \textbf{Domain-specific semantics}: IaC faults span syntactic errors, cross-resource dependency issues, and cloud provider-specific constraints that differ fundamentally from traditional programming bugs.
    \item \textbf{Validation complexity}: Correctness requires interaction with cloud provider APIs and validation tools (e.g., \texttt{terraform validate}, \texttt{cfn-lint}), not just syntactic well-formedness.
    \item \textbf{Lack of benchmarks}: No standardized benchmark exists for evaluating LLM-based IaC repair, making it difficult to compare approaches or measure progress.
\end{itemize}

We address these gaps with \textbf{Cloud-Gym}, a framework that provides:
\begin{enumerate}
    \item A \textbf{37-fault taxonomy} across 8 categories, covering both Terraform/HCL and CloudFormation/YAML fault types.
    \item A \textbf{programmatic inversion engine} that generates validated broken configurations from gold (valid) instances through targeted fault injection.
    \item A \textbf{training data pipeline} producing thousands of (broken, gold, error) triples for supervised fine-tuning.
    \item An \textbf{evaluation harness} with pass@$k$ metrics and stratified reporting by fault category, difficulty, and IaC format.
\end{enumerate}

% TODO: Add contributions summary and paper outline

\section{Related Work}
\label{sec:related_work}

\paragraph{IaC Analysis and Repair.}
Prior work on IaC analysis has focused on static analysis tools~\citep{cfnlint2023}, security scanning~\citep{checkov2023}, and policy enforcement. However, automated \emph{repair} of IaC configurations using learned models remains largely unexplored. Existing tools can detect issues but require manual intervention for fixes.

\paragraph{LLM-Based Code Repair.}
The application of LLMs to automated program repair has gained significant attention. Benchmarks like SWE-bench~\citep{jimenez2024swebench} evaluate end-to-end issue resolution, while HumanEval~\citep{chen2021codex} and MBPP~\citep{austin2021program} focus on code generation. CLI-Gym~\citep{cligym2024} provides a reinforcement learning environment for command-line tasks. Cloud-Gym extends this line of work to the IaC domain.

\paragraph{Fault Injection and Mutation Testing.}
Mutation testing~\citep{jia2011mutation} systematically introduces faults to evaluate test suite adequacy. Cloud-Gym's inversion engine draws on mutation testing principles but targets IaC-specific fault patterns validated by real infrastructure tools.

% TODO: Expand with additional citations once literature review is complete

\section{Method}
\label{sec:method}

\subsection{Fault Taxonomy}
\label{sec:taxonomy}

We define a taxonomy of 37 fault types organized into 8 categories, designed to cover the spectrum of real-world IaC misconfigurations. Table~\ref{tab:taxonomy} summarizes the taxonomy.

% TODO: Insert taxonomy table (auto-generated from code)
% \input{figures/taxonomy_table}

\begin{itemize}
    \item \textbf{Syntactic} (7 faults): Parse-level errors including missing braces, type mismatches, invalid syntax, and missing required arguments.
    \item \textbf{Reference} (6 faults): Broken variable references, non-existent resource references, undefined parameters, and invalid module sources.
    \item \textbf{Semantic} (7 faults): Logically invalid values such as wrong resource types, malformed AMI IDs, invalid CIDR blocks, and non-existent regions.
    \item \textbf{Dependency} (4 faults): Circular dependencies and missing ordering constraints.
    \item \textbf{Provider} (2 faults): Missing provider configurations and impossible version constraints.
    \item \textbf{Security} (4 faults): Overly permissive network rules and missing encryption.
    \item \textbf{Cross-Resource} (3 faults): Mismatched VPC/subnet references and wrong availability zones.
    \item \textbf{Intrinsic} (4 faults, CF-only): Malformed \texttt{!Sub}, out-of-bounds \texttt{!Select}, undefined conditions, and invalid \texttt{!Join}.
\end{itemize}

\subsection{Inversion Engine}
\label{sec:inversion}

The inversion engine transforms gold (valid) configurations into broken ones via targeted fault injection. Given a gold config $c_g$ and fault type $f$, the engine produces a broken config $c_b$ such that:
\begin{enumerate}
    \item $c_b \neq c_g$ (the config was actually modified),
    \item $\text{validate}(c_b) = \text{FAIL}$ (the modification breaks validation), and
    \item The diff between $c_g$ and $c_b$ is minimal (single-fault injection).
\end{enumerate}

\paragraph{Programmatic Injection.} For Terraform/HCL, we use a parse-then-regex approach: \texttt{python-hcl2} parses the config for structural analysis, then targeted string manipulation applies the fault. For CloudFormation, we use YAML/JSON round-trip mutation. Each fault type has a dedicated injector function.

\paragraph{Agentic Injection.} We also employ local LLMs (via Ollama) for diverse fault generation, with quality gates ensuring similarity ($\geq 0.7$) and minimal diff ($<20\%$ lines changed).

\subsection{Evaluation Protocol}
\label{sec:evaluation}

We evaluate repair models using the pass@$k$ metric~\citep{chen2021codex}:
\[
\text{pass@}k = \mathbb{E}\left[1 - \frac{\binom{n-c}{k}}{\binom{n}{k}}\right]
\]
where $n$ is the number of attempts and $c$ is the number of correct repairs. A repair is correct if the repaired config passes validation with zero errors.

% TODO: Add pipeline diagram figure

\section{Dataset}
\label{sec:dataset}

\subsection{Gold Configuration Collection}
\label{sec:gold_collection}

We collect gold (valid) IaC configurations from three sources:
\begin{itemize}
    \item \textbf{GitHub}: Public repositories containing Terraform and CloudFormation files, filtered by minimum resource count and validated.
    \item \textbf{Terraform Registry}: Official modules from the Terraform Registry.
    \item \textbf{AWS Samples}: CloudFormation templates from AWS's official sample repositories.
\end{itemize}

All collected configurations are validated using \texttt{terraform validate} (for Terraform) or \texttt{cfn-lint} (for CloudFormation). Only configurations that pass validation are retained as gold instances.

% TODO: Add table with collection statistics

\subsection{Training Data Generation}
\label{sec:training_data}

From the gold configurations, we generate training pairs using the inversion engine (Section~\ref{sec:inversion}). For each gold config, we apply $N$ programmatic fault injections via stratified sampling across categories, plus $M$ agentic injections for diversity.

The resulting dataset is split by gold config hash to prevent leakage: all variants derived from the same gold config appear in the same split (train/val/test at 80/10/10).

% TODO: Add dataset statistics table

\subsection{Benchmark Curation}
\label{sec:benchmark}

From the test split, we curate a balanced benchmark subset with the following criteria:
\begin{itemize}
    \item Single-fault configurations only
    \item Minimum 10-line configs (non-trivial)
    \item Balanced across fault categories and difficulty levels
    \item Deduplicated per gold config
\end{itemize}

% TODO: Add benchmark statistics table

\section{Experiments}
\label{sec:experiments}

\subsection{Baseline Models}
\label{sec:baselines}

We evaluate two open-weight models in a zero-shot repair setting:
\begin{itemize}
    \item \textbf{DeepSeek-R1 1.5B}: A small reasoning model, serving as a performance floor.
    \item \textbf{Qwen2.5-Coder 7B}: A code-specialized model, serving as the main baseline.
\end{itemize}

Both models are run locally via Ollama. The repair prompt provides the broken configuration and its validation errors, asking the model to return only the fixed configuration.

\subsection{Results}
\label{sec:results}

% TODO: Insert results tables (auto-generated from baseline runs)

% \begin{table}[h]
% \centering
% \caption{Overall pass@$k$ results.}
% \label{tab:overall_results}
% \begin{tabular}{lcc}
% \toprule
% Model & pass@1 & pass@3 \\
% \midrule
% DeepSeek-R1 1.5B & -- & -- \\
% Qwen2.5-Coder 7B & -- & -- \\
% \bottomrule
% \end{tabular}
% \end{table}

\subsection{Analysis}
\label{sec:analysis}

% TODO: Add analysis by category, difficulty, and format
% - Syntactic faults are easiest to repair
% - Semantic and cross-resource faults are most challenging
% - CF vs TF comparison
% - Difficulty level correlation with repair success

\section{Conclusion}
\label{sec:conclusion}

We presented Cloud-Gym, a benchmark and data generation framework for evaluating LLM-based repair of Infrastructure-as-Code configurations. Our 37-fault taxonomy covers the spectrum of real-world IaC misconfigurations across Terraform and CloudFormation. The programmatic inversion engine generates validated training pairs at scale, and the evaluation harness provides fine-grained pass@$k$ metrics.

Our baseline results indicate that current open-weight models show varying repair capability across fault categories, with syntactic faults being most amenable to repair and semantic/cross-resource faults posing the greatest challenge. This gap motivates future work in IaC-specific model training and more sophisticated repair strategies.

\paragraph{Future Work.}
\begin{itemize}
    \item Fine-tuning models on Cloud-Gym training data
    \item Extending the taxonomy to additional IaC tools (Pulumi, CDK, Ansible)
    \item Multi-fault repair scenarios
    \item Integration with CI/CD pipelines for real-time repair suggestions
\end{itemize}

% TODO: Finalize based on experimental results


\bibliographystyle{plainnat}
\bibliography{references}

\end{document}
